\begin{proof}
The emission distribution at latent state $j$ is $C_{\cdot j}$.
By Lemma~\ref{II:M:lem:projective-column-limit}, its KL divergence from
$\widehat C_{\cdot j}$ is at most $\Psi(\eta_{\rm obs})$.
Given $(j,a)$, the next-state distributions are
\[
 R_{a,j}=\frac{T_ae_j}{C_{aj}},\qquad
 \widehat R_{a,j}=\frac{\widehat T_ae_j}{\widehat C_{aj}}.
\]
The same lemma bounds their KL divergence by $\Psi(\eta)$ and
ensures matching support. The joint one-step kernels $K_j(a,i)$ and
$\widehat K_j(a,i)$ therefore satisfy, by the finite conditional
probability identity,
\[
 D_{\rm KL}(K_j\|\widehat K_j)
 =D_{\rm KL}(C_{\cdot j}\|\widehat C_{\cdot j})
  +\sum_a C_{aj}D_{\rm KL}(R_{a,j}\|\widehat R_{a,j})
 \le\Psi(\eta_{\rm obs})+\Psi(\eta).
\]
Expanding the log-likelihood ratio of the joint latent-state and letter
trajectory gives its initial KL plus the sum of the conditional kernel
KL divergences. The initial KL is at most $\Psi(e_0)$, so the displayed
bound follows for the joint trajectory. Marginalizing the latent states
cannot increase KL: grouping terms with the same observed word and
applying the log-sum inequality proves this directly. Finally apply
$\Psi(e)\le e^2/8$ to all three terms.
\end{proof}
